\chapter[Train Track Case Analysis: The Candelabra]{Train Track Case Analysis: The Candelabra}
\label{app:case analysis candelabra}

First notice that the triangle must be rotated by $f_\beta$, otherwise the image of all of the real edges will trace. There are two possible rotations of the triangle, one is the result of a single clockwise rotation, and the other is the result of two clockwise rotations. 

First we assume that the triangle is rotated once clockwise. Thus, $\fbi$ begins with $r^*$, but could only be $r^\circ$ or $r^+$, and $\fbp$ begins with $b^*$. On the other hand, $\fbr$ and $\fbb$ begin with $o^*$ or $s^*$, while $\fbo$ and $\fbs$ begin with $i$ or $p$. As shown in ????

\subsection{NHT}\label{NHT}
Suppose
\[
\fbi = r^\circ
\]
and
\[
\fbp = b^\circ
\]

Suppose further that 

\[
\fbr = o^\circ
\]
and
\[
\fbs = i^\circ
\]

This is NHT1.

This then forces $\fbb$ to begin with either $s^\circ$ or $s^-$, so suppose it were the former. Then 
\[
\fbo = i^\circ
\]
is forced.

This is NHT2. This is Candidate 1.

So suppose now
\[
\fbb = s^-...
\]
This is NHT3.

Here $\fbb$ must continue with $i^-$, to give us
\[
\fbb = s^- i^- s^*...
\]

But now $\fbo$ will be absorbed:
\[
\fbo = i^+ s^+ \Abso{\fbr}{\fbb}
\]

shown in NHT4.

So now suppose 
\[
\fbs = p^+...
\]
The next letter must be $o^*$ or it will trace, but now both $\fbb$ and $\fbs$ will be absorbed, shown in NHT5.

So now assume
\[
\fbr = o^-...
\]

This is NHT6.

Here $\fbr$ can continue with $i$ or $p$.

\subsubsection{Jehu}\label{Jehu}
Suppose
\[
\fbr = o^- i^*...
\]

This is Jehu1.

Here the next letter can be $i^\circ$ or $i^+$. Suppose it were the former, then $\fbo$ and $\fbs$ are forced to begin with $p^*$, in particular $\fbs$ must begin with $s^+...$

This is Jehu2.

But here $\fbb$ is forced to begin with $s^-$ and to the right of the incoming $\fbs$, then it continues with $i^+ s^*...$, or it will be absorbed. Then Jehu3 is forced.

Here $\fbo$ must begin and end at $p^\circ$ otherwise it will immediately trace. Here $\fbs$ must end at $o^\circ$, or it will either trace or be absorbed. This forces $\fbb$ to also end. This gives us Jehu4. This is a candidate train track map, but in fact it will be subsumed in Candidate Family 2, to follow.

Now suppose
\[
\fbr = o^- i^+...
\]

Then the situation in Jehu5 is forced. Here $\fbr$ and $\fbb$ twist counterclockwise between $s$ and $i$ and end there. This gives Candidate Family 2.

\subsubsection{Kexy}
Back to the end of NHT \ref{NHT}. Suppose
\[
\fbr = o^- p^*...
\]

Then we have Kexy1.

Here $\fbr$ must continue with $r^-$ and to the right of both $\fbo$ and $\fbs$. Otherwise, at least one of $\fbo$ or $\fbs$ will trace. But now it traces immediately:
\[
\fbr = o^- i^- r^\times
\]

This finishes the situation under NHT \ref{NHT}.

\subsection{CMJ}\label{CMJ}
Suppose
\[
\fbi = r^+...
\]
and
\[
\fbp = b^\circ
\]

This is CMJ1.

Here $\fbi$ can continue with $p^\circ$ or $p^+$.

\subsubsection{Mome}\label{Mome}
Suppose
\[
\fbi = r^+ p^\circ
\]

This is the situation depicted in Mome1.

Here $\fbr$ must begin with $s$ or $o$. We claim that the former cannot occur. Indeed, if $\fbr = s^*....$ then $\fbb$ also starts with $s$. But then $\fbr$ must in fact start with $s^-$ and to the right of $\fbb$ otherwise $\fbb$ will trace, but then we have Mome2. Here if $\fbr$ continued with $i^\circ$ or $i^+$ then $\fbo$ and $\fbs$ will be absorbed between $\fbp$ and $\fbi$; but then if $\fbr$ continued with $i^-$ then it immediately traces.

Thus we must have 
\[
\fbr = o^\circ
\]
or
\[
\fbr = o^-...
\]

\subsubsubsection{Kedge}\label{Kedge}
Suppose
\[
\fbr = o^\circ
\]

Then we must have
\[
\fbb = s^*...
\]

This is Kedge1.

Here $\fbb$ either ends at $s^\circ$ or continues with $s^-$.

\subsubsubsubsection{Dindle}\label{Dindle}
Suppose
\[
\fbb = s^\circ
\]

This is Dindle1.

Here $\fbs$ can either end at $i^\circ$ or continue with $i^-$. Suppose
\[
\fbs = i^\circ
\]
Then we have Dindle2, which is forced. This is a candidate train track map that will be subsumed in Candidate Family 3, to follow.

Suppose instead
\[
\fbs = i^-...
\]

Then we have Dindle3. Here $\fbo$ and $\fbs$ twist clockwise between $r$ and $i$, and end there. This gives Candidate Family 3.

\subsubsubsubsection{Jirble}
Back to the end of Kedge \ref{Kedge}. Suppose
\[
\fbb = i^-...
\]

Then we have Jirble1.

Here both $\fbo$ and $\fbb$ must continue with $i^-$ and to the right of $\fbs$. Otherwise, $\fbs$ will be forced to begin with $i^+$ and immediately trace.

Thus we have
\[
\fbo = i^-...
\]
and 
\[
\fbb = s^- i^-...
\]

This is Jirble2.

Here $\fbb$ must continue with $r^-$, otherwise $\fbo$ will be forced to continue with $r^+$ and to the left of $\fbb$. In which case it will be absorbed:
\[
\fbo = i^- \hat{r}^+ p^- r^- \Abso{\fbp}{\fbi}
\]
or
\[
\fbo = i^- \hat{r}^+ i^+ s^+ \Abso{\fbr}{\fbb}
\]

Thus we must have
\[
\fbb = s^- i^- \hat{r}^- i^*...
\]
which forces
\[
\fbs = i^- r^*...
\]

This is Jirble3.

Now notice that $\fbb$, $\fbo$, and $\fbs$ are stuck in a three-way clockwise path mill between $i$ and $r$.

This is a contradiction.


\subsubsubsection{Mulse}
Back to the end of Mome \ref{Mome}. Suppose
\[
\fbr = o^-...
\]

This is Mulse1.

Here if $\fbr$ were to end with $i^\circ$ or continue with $i^+$, then $\fbo$ and $\fbs$ would both begin with $i^+$, but then either
\[
\fbs = i^+ s^\times
\]
or
\[
\fbo = i^+ o^\times
\]
would occur. Thus we must have $\fbr$ continue with $i^-$ and to the right of $\fbo$ and $\fbs$. But then it traces:
\[
\fbr = o^- i^- r^\times
\]

This is a contradiction.

\subsubsection{Nesh}\label{Nesh}
Back to the end of CMJ \ref{CMJ}. Suppose
\[
\fbi = r^+ p^+...
\]

This is Nesh1.

Here $\fbi$ can continue with $r^*$, $s^*$, or $o^*$.

\subsubsubsection{Nudzh}\label{Nudzh}
Suppose $\fbi$ continued with $r^*$. This forces
\[
\fbi = i^*...
\]
and
\[
\fbo = i^*...
\]

where for $\fbo$ the first letter could either be $i^\circ$ or $i^-$ (it cannot be $i^+$ as either itself would trace or $\fbs$ would).

Suppose 
\[
\fbo = i^\circ
\]

This forces Nudzh2. 

Here $\fbr$ and $\fbb$ must begin with $o^-$ and to the right of $\fbs$. Otherwise, $\fbs$ would be absorbed:
\[
\fbs = i^+ \hat{o}^+ b^+ \Abso{\fbp}{\fbi}
\]

But now $\fbr$ will trace:
\[
\fbr = o^- i^- r^\times
\]

Back to the situation depicted in Nudzh2. Suppose
\[
\fbo = i^-...
\]

This is Nudzh3. 

Here we claim that $\fbr$ must begin and end at $o^\circ$. Indeed, if it began with $o^+$ then it will be absorbed:
\[
\fbr = o^+ b^+ \Abso{\fbp}{\fbi}
\]

If it began with $o^-$ it would trace;:
\[
\fbr = o^- i6- r^\times
\]

If it began with $s^*$, then $\fbb$ would also begin with $s^*$. But then $\fbb$ must begin with $s^+$ and to the left of $\fbr$ otherwise $\fbr$ would again trace. But then we have
\[
\fbb = s^+ b^\times
\]

Therefore, we must have
\[
\fbr = o^\circ
\]
which forces $\fbb = s^\circ$ or $\fbb = s^-...$

This is Nudzh4.

Here $\fbb$ can either begin and end at $s^\circ$, or begin with $s^-$.

\subsubsubsubsection{Keffel}\label{Keffel}
Suppose
\[
\fbb = s^\circ
\]

This is Keffel1.

Here $\fbs$ can begin and end at $i^\circ$, or begin with $i^-$. In the former case, $\fbi$ and $\fbo$ then twist between $r$ and $p$, giving Candidate Family 4. This is Keffel2.

Suppose it were the latter case:
\[
\fbs = i^-...
\]

Then we have Keffel3.

Here $\fbi$ twist between $r$ and $p$ and end there as it cannot escape or it will trace. Then $\fbs$ and $\fbo$ twist clockwise between $i$ and $r$/$p$ (depending on which is left), and end there as well. This gives Candidate Family 5. An example is shown in Keffel4.


\subsubsubsubsection{Mizzle}
Back to the end of Nudzh\ref{Nudzh}. Suppose
\[
\fbb = s^-...
\]
and it continues:
\[
\fbb = s^- i^- r^*...
\]

This is Mizzle1.

Here $\fbs$ can either begin and end at $i^\circ$, or begin with $i^-...$

Suppose it were the former. Then we have Mizzel2. Here both $\fbb$ and $\fbo$ must continue with $r^+$ and to the left of $\fbi$, otherwise $\fbi$ will immediately trace. This gives Mizzle3. Here $\fbb$, $\fbo$, and $\fbi$ are in a three-way counterclockwise path mill between $p$ and $r$. This is a contradiction.

Suppose it were the latter case:
\[
\fbs = i^-...
\]

This forces the situation depicted in Mizzle4. 

Here $\fbb$, $\fbo$, and $\fbs$ all have to continue with $r^+$ and to the left of $\fbi$, otherwise $\fbi$ will be forced to continue with $r^-$ and immediately trace,

This is Mizzle5.

Here $\fbi$ must eventually end at $r$ or $p$, possible after twisting between the two. But then $\fbb$, $\fbo$, and $\fbs$ will be stuck in a three-way clockwise path mill between which ever is left ($r$ or $p$), and $i$. 

This is a contradiction.

\subsubsubsection{Snirt}\label{Snirt}
Back to the end of Nesh \ref{Nesh}. Suppose $\fbi$ continues with $o^*$:
\[
\fbi = r^+ p^+ \hat{o}^*...
\]

This is Snirt1.

Here $\fbi$ can either continue with $o^\circ$ or $o^+$.

\subsubsubsubsection{Tittup}\label{Tittup}
Suppose $\fbi$ continued with $o^\circ$:
\[
\fbi = r^+ p^+ \hat{o}^\circ
\]

This is Tittup1.

Here if $\fbs$ were to begin with $i^+$, then we have
\[
\fbs = i^+ o^- i^*...
\]
Here the third letter cannot be $p$ as in that case it would be absorbed:
\[
\fbs = i^+ o^- \hat{p}^- r^- \Abso{\fbp}{\fbi}
\]

Thus we have Tittup2.

Here $\fbs$ must continue with $i^+$ and to the left of both $\fbr$ and $\fbb$. Otherwise, $\fbr$ would be forced to continue with $i^-$, in which case it would trace:
\[
\fbr = o^- \hat{i}^- o^+ p^- r^\times
\]

Thus we must have
\[
\fbs = i^+ o^- \hat{i}^+...
\]
but now it either traces:
\[
\fbs = i^+ o^- \hat{i}^+ s^\times
\]
or is absorbed:
\[
\fbs = i^+ o^- \hat{i}^+ o^+ b^+ \Abso{\fbp}{\fbi}
\]

Therefore, in the situation depicted in Tittup1, $\fbo$ must begin with $p^-$ and to the right of $\fbs$:
\[
\fbo = p^- r^*...
\]

This gives Tittup3. 

Here $\fbo$ can either end at $r^\circ$ or continue with $r^-$.

\subsubsubsubsubsection{Epigone}\label{Epigone}
Suppose $\fbo$ ends at $r^\circ$:
\[
\fbo = p^- r^\circ
\]

This forces
\[
\fbs = p^\circ
\]

as otherwise $\fbs$ wither traces:
\[
\fbs = p^- r^+ p^+ o^- i^+ s^\times
\]
or it is absorbed:
\[
\fbs = p^- r^+ p^+ o^- p^- r^- \Abso{\fbp}{\fbi}
\]

Thus we must have Epigone1.

Here $\fbb$ and $\fbr$ twist between $s$ and $i$ $k$ times and end there, giving us Candidate Family 6.

\subsubsubsubsubsection{Farrago}\label{Farrago}
Back to the end of Tittup \ref{Tittup}. Suppose $\fbo$ were to continue with $r^-$:
\[
\fbo = p^- r^- p^*...
\]
which forces:
\[
\fbs = p^- r^*...
\]

This is Farrago1.

Here $\fbo$ and $\fbs$ must end at $r$ or $p$ after possibly twisting among themselves $m$ times. The two cannot escape $p$ and $r$ as either they would trace or be absorbed. On the other hand $\fbb$ and $\fbr$ also twist $k$ times between $s$ and $i$ and end there. This gives us Candidate Family 7.

\subsubsubsubsection{Ullage}
Back to the end of Snirt. Suppose $\fbi$ continued with $o^+$:
\[
\fbi = r^+ p^+ \hat{o}^+ p^*...
\]

This is Ullage1.

Here $\fbi$ must continue with $p^-$ and to the right of $\fbo$ and $\fbs$. Otherwise, $\fbo$ will be forced to begin with $p^+$, after which it will immediately trace.

Thus we must have
\[
\fbi = r^+ p^+ o^+ \hat{p}^- r^*...
\]

This is Ullage2.

Here $\fbs$ can either begin and end at $i^\circ$, or begin with $i^-$, or $i^+$.

Suppose
\[
\fbs = p^\circ
\]

Then this forces
\[
\fbo = p^- r^*...
\]

This is the situation depicted in Ullage3. 

Here if $\fbi$ were to continue with $r^\circ$ or $r^+$, then $\fbo$ is forced to continue with $r^+$, but then it will trace:
\[
\fbo = p^- \hat{r}^+ p^+ o^\times
\]

Thus $\fbi$ must continue with $r^-$ and to the right of $\fbo$. But then it will be absorbed:
\[
\fbi = ...r^- p^+ \Abso{\fbo}{\fbs}
\]

This is a contradiction.

So suppose instead
\[
\fbs = p^+...
\]
then it continues:
\[
\fbs = p^+ o^- i^*...
\]

This is Ullage4.

Here if $\fbb$ were to continue with $i^\circ$ or $i^+$ then $\fbs$ must continue with $i^+$, in which case it either immediately traces or is absorbed between $\fbp$ and $\fbi$. Hence, $\fbb$ and $\fbr$ must both continue with $i^-$ and to the right of $\fbs$. But now $r$ will trace:
\[
\fbr = o^- \hat{i}^- o^+ p^- r^\times
\]

This is a contradiction.

Thus it must be the case that
\[
\fbs = p^- r^*...
\]
which forces
\[
\fbo = p^- r^*...
\]

This is Ullage5. 

Here $\fbi$ must continue with $r^-$ and to the right of $\fbo$ and $\fbs$. Otherwise, $\fbo$ will be forced to continue with $r^+$, in which case it will trace:
\[
\fbo = p^- \hat{r}^+ p^+ o^\times
\]

This gives Ullage6.

Here $\fbi$, $\fbo$, and $\fbs$ are in a three-way clockwise path mill between $p$ and $r$.

This is a contradiction.


\subsubsubsection{Quire}\label{Quire}
Back to the end of Nesh \ref{Nesh}. Suppose $\fbi$ continued with $s^*$:
\[
\fbi = r^+ p^+ s^*...
\]

This is Quire1.

Here $\fbr$ can begin with $s^*$, or $o^*$. 

\subsubsubsubsection{Yepsen}\label{Yepsen}
Suppose 
\[
\fbr = s^*...
\]
Then in fact must be $s^-$ and to the right of $\fbb$, otherwise $\fbb$ would be absorbed. After $s^-$, the next letter must be $i$. Indeed, otherwise it either traces:
\[
\fbr = s^- p^- r^\times
\]
or it would force $\fbi$ to be absorbed:
\[
\fbi = r^+ p^+ o^+ b^+ \Abso{\fbp}{\fbi}
\]

This gives Yepsen1.

This then forces 
\[
\fbb = s^- i^*...
\]

This is Yepsen2.

Here $\fbi$ can either continue with $s^\circ$, or $s^+$. Whereas $s^-$ is not possible as that would mean it either immediately traces, or eventually absorbed:
\[
\fbi = ...s^- p^- r^- \Abso{\fbp}{\fbi}
\]

\subsubsubsubsection{Gubbins}
Suppose $\fbi$ ends at $s^\circ$:
\[
\fbi = r^+ p^+ s^\circ
\]

This forces Gubbins1.

Here $\fbb$ cannot continue with $i^\circ$ or $i^-$, otherwise $\fbr$ will be forced to continue with $i^-$. In which case $\fbr$ will trace:
\[
\fbr = s^- \hat{i}^- s^+ p^- r^\times
\]

Thus we must have
\[
\fbb = s^- \hat{i}^+...
\]
but then it will trace:
\[
\fbb = s^- \hat{i}^+ s^+ b^\times
\]

shown in Gubbins2.

This is a contradiction.

\subsubsubsubsection{Ignavia}
Back to the end of Yepsen\ref{Yepsen}. Suppose $\fbi$ continues with $s^+$:
\[
\fbi = r^+ p^+ \hat{s}^+ p^*...
\]
This is Ignavia1.

Here again $\fbr$ is forced to continue with $i^-$ otherwise $\fbb$ will be absorbed. 

This gives Ignavia2.

Here $\fbr$ must continue with $p^-$ and to the right of $\fbs$. Otherwise, $\fbs$ will be forced to continue with $p^+$ and immediately trace. But then $\fbr$ traces:
\[
\fbr = s^- i^- s^+ \hat{p}^- r^\times
\]

This is a contradiction.

\subsubsubsubsection{Zoilus}\label{Zoilus}
Back to the end of Quire \ref{Quire}. Suppose
\[
\fbr = o^*...
\]

This is Zoilus1.

Here $\fbi$ can continue with $s^\circ$ or $s^+$.

\subsubsubsubsubsection{Jumbuck}\label{Jumbuck}
Suppose $\fbi$ continued with $s^\circ$:
\[
\fbi = r^+ p^+ s^\circ
\]

This is Jumbuck1.

Here if $\fbs$ were to begin with $p^+$, then we would have Jumbuck2. But here $\fbr$ must begin with $o^-$ and to the right of $\fbo$ to prevent $\fbo$ from being absorbed. But then $\fbr$ will trace:
\[
\fbr = o^- p^- r^\times
\]

Hence, $\fbs$ must begin with $p^\circ$ or $p^-$, which means $\fbo$ must begin with $p^-$. This forces 
\[
\fbr = o^\circ
\]
as otherwise $\fbr$ will trace. This in turn forces
\[
\fbb = s^- i^*....
\]

This is Jumbuck3.

Here $\fbb$ can either end at $i^\circ$ or continue with $i^-$.

Suppose it ended at $i^\circ$:
\[
\fbb = s^- i^\circ
\]
Then we have the situation depicted in Jumbuck4, where $\fbs$ and $\fbo$ twist $k$ times between $p$ and $r$ and end there, giving Candidate Family 8.

Now suppose instead
\[
\fbb = s^- i^-...
\]
then it continues:
\[
\fbb = s^- i^- s^+ p^- r^*...
\]

This is Jumbuck5. 

Here $\fbs$, $\fbb$, and $\fbo$ are stuck in a three-way clockwise path mill between $p$ and $r$. 

This is a contradiction.

\subsubsubsubsubsection{Nithing}\label{Nithing}
Back to the end of Zoilus \ref{Zoilus}. Suppose $\fbi$ were to continue with $s^+$. Then it continues:
\[
\fbi = r^+ p^+ \hat{s}^+ p^- r^*...
\]

This is Nithing1.

Here as before we must have $\fbo$ begin with $p^-$. Otherwise, $\fbs$ will begin with $p^+$, in which case it will then force $\fbr$ to begin with $o^-$ and then trace.

Thus we have the situation depicted in Nithing2.

Here $\fbs$ must end at either $p$ or $r$, possibly after twisting $m$ times between the two. It can never escape or it will trace. Once it ends, $\fbi$ must then end at either $s$ or $p$/$r$, possibly after twisting $k$ times. Once this is decided, $\fbb$ and $\fbo$ must immediately end without more twisting, as otherwise one of them will eventually escape to the outer perimeter and therefore be absorbed. An example of such a situation is shown in Nithing3, where one can see $\fbb$ will eventually escape to the outer perimeter. 


This gives Candidate Family 9.

This finishes the analysis under the assumptions of CMJ \ref{CMJ}.

\subsection{WST}\label{WST}
Suppose
\[
\fbi = r^\circ
\]
and
\[
\fbp = b^-...
\]

This is WST1.

here $\fbp$ can continue with $o^*$ or $s^*$.

\subsubsection{Whid}\label{Whid}
Suppose $\fbp$ continued with $o^*$. Then it is either $o^\circ$ or $o^-$.

\subsubsubsection{Tilly}
Suppose 
\[
\fbp = b^- o^\circ
\]

This forces 
\[
\fbb = s^*...
\]
and
\[
\fbr = s^-...
\]

This is Tilly1.

Here the next letter for $\fbr$ is either $i^*$ or $p^*$.

Suppose it were the former:
\[
\fbr = s^- i^...
\]

Then both $\fbs$ and $\fbo$ must begin at $p^*$. Indeed, otherwise, $\fbs$ would be forced to begin with $p^+$, which it would trace; or $\fbr$ would continue with $p^-$, in which case it would trace. 


This is Tilly2.

Then $\fbs$ must in fact begin with $p^+$, otherwise $\fbo$ would be absorbed. But then $\fbs$ immediately traces.

So suppose instead $\fbr$ continued with $p^*$. This is Tilly3.

Here if $\fbr$ were to continue with $p^\circ$ or $p^+$, then $\fbs$ will begin with $p^+$, in which case it will immediately trace. But this means $\fbr$ must continue with $p^-$, in which case it immediately traces.

This is a contradiction.


\subsubsubsection{Whelm}\label{Whelm}
Back to the end of Whid \ref{Whid}. Suppose
\[
\fbp = b^- o^-...
\]

This is Whelm1.

Here $\fbo$ can either begin with $i^*$ or $p^*$.

\subsubsubsubsection{Bodger}
Suppose
\[
\fbo = i^*...
\]

This forces the situation depicted in Bodger1.

Here $\fbo$ must end at $s$ or $i$, possibly after some twisting between the two. Then $\fbb$ will end at $o$ or $s$/$i$ possibly after twisting. But now at least one of $\fbp$ or $\fbr$ will be forced to escape to the outer perimeter, thereby eventually absorbed. An example is shown in Bodger2. Here $\fbr$ must continue with $s^+$, or $|fbp$ will be absorbed between $\fbr$ and $\fbb$. But now $\fbr$ will escape to the outer perimeter:
\[
\fbr = o^- i^+ \hat{s}^+ i^- o^+...
\]
or
\[
\fbr = o^- i^+ \hat{s}^+ i^- s^-...
\]

\subsubsubsubsection{Callid}
Back to the end of Whelm \ref{Whelm}. Suppose
\[
\fbo = p^*...
\]

Then it must in fact be
\[
\fbo = p^\circ
\]

Otherwise it would immediately trace or be absorbed.

Then this forces
\[
\fbs = p^+ o^*...
\]
which forces
\[
\fbb = o^-...
\]

This is Callid2.

Here if $\fbs$ were to end at $o^\circ$, then $\fbp$, $\fbr$, and $\fbb$ will be stuck in a three-way counterclockwise path mill between $i$ and $s$.

So suppose $\fbs$ here continued with $o^-...$

This is Callid3.

Here $\fbp$, $\fbr$, and $\fbb$ must contiue with $i^-$ and to the right of $\fbs$. Otherwise, $\fbs$ will trace. 

This is Callid4.

Here $\fbs$ must eventually end at $i$ or $o$ after possibly twisting between the two. But then $\fbp$, $\fbb$, and $\fbr$ will be stuck in a three-way path mill.

This is a contradiction.

\subsubsection{Dido}
Back to the end of WST \ref{WST}. Suppose $\fbp$ were to continue with $s^*$. 

This is Dido1.

Here $\fbr$ and $\fbp$ must continue with $s^-$ and to the right of $\fbb$. Otherwise, $\fbb$ will trace. Here $\fbr$ can either continue with $p^*$ or $i^*$. If it were the former, then we have Dido2. Here if $\fbr$ were to continue with $p^\circ$ or $p^+$ then $\fbs$ will immediately trace; but then if $\fbr$ continued with $p^-$ it will immediately trace. Thus $\fbr$ must continue with $i^*$. This is Dido3.

Here both $\fbo$ and $\fbs$ must begin with $p^*$, otherwise either $\fbp$ will be absorbed:
\[
\fbp = b^- s^- \hat{i}^- r^- \Abso{\fbp}{\fbi}
\]
or $\fbs$ will trace:
\[
\fbs = i^+ s^\times
\]

But now if $\fbs$ began with $p^\circ$ or $p^-$, then $\fbo$ will begin with $p^-$ and be absrobed:
\[
\fbo = p^0 r^- \Abso{\fbp}{\fbi}
\]

But then this means $\fbs$ begins with $p^+$, which mean it will immediately trace:
\[
\fbs = p^+ s^\times
\]

This is a contradiction.

This completes the analysis of the scanario under the assumptions of WST \ref{WST}.

\subsection{PPG}
Suppose
\[
\fbi = r^+...
\]
and
\[
\fbp = b^-...
\]

This is PPG1.

Here $\fbp$ can continue with $o^*$ or $s^*$.

\subsubsection{Fash}\label{Fash}
Suppose $\fbp$ continued with $o^*$. It can either be $o^\circ$ or $o^-$.

On the other hand, $\fbi$ can either continue with $p^\circ$ or $p^+$.

Suppose
\[
\fbp = b^- o^\circ
\]
and
\[
\fbi = r^+ p^\circ
\]

This is Fash1.

Here $\fbo$ must begin with $i^-$ and to the right of $\fbs$. Otherwise, $\fbs$ will be forced to begin with $i^+$ in which case it immediately traces. 

Now we have Fash2.

Now if $\fbr$ began with $s^\circ$ or $s^+$, $\fbb$ will begin with $s^+$ and immediately trace. But then if $\fbr$ began with $s^-$, then it also traces;
\[
\fbr = s^- i^- r^\times
\]

So suppose instead 
\[
\fbp = b^- o^-...
\]
and
\[
\fbi = r^+ p^\circ
\]

This is Fash3.

Here if $\fbp$ were to continue with $i^*$, then we will have the situation depicted in Fash4.

Here $\fbs$ must begin with $i^\circ$ or $i^-$, or it will immediately trace. But this will force $\fbr$ to continue with $i^-$ and immediately trace.

So suppose instead $\fbp$ were to continue with $b^*$. This gives Fash5.

Here again $\fbs$ must begin with $i^\circ$ or $i^-$, but again $\fbr$ will trace.


So suppose instead 
\[
\fbp = b^- o^-...
\]
and
\[
\fbi = r^+ p^+...
\]

This is Fash6.

Here $\fbi$ can continue with $o^*$ or $s^*$.

\subsubsubsection{Xeric}\label{Xeric}
Suppose $\fbi$ continued with $o^*$:
\[
\fbi = r^+ p^+ o^*...
\]

This forces the situation depicted in Xeric1.

Here $\fbr$ can continue into 1. or 2.

Suppose it continued into 1. Then we have Xeric2.

Here if $\fbr$ were to continue with $p^\circ$ or $p^+$, then $\fbo$ will be forced to begin with $p^+$ and it will immediately trace. So then $\fbr$ must continue with $p^-$ but then it immediately traces. Thus $\fbr$ cannot go into 1.

Suppose it continued into 2. Then it forces
\[
\fbb = o^-...
\]

This is the situation depicted in Xeric3.

Here $\fbi$ can either continue with $o^\circ$ or $o^+$.

\subsubsubsubsection{Fizgig}\label{Fizgig}
Suppose $\fbi$ were to continue with $o^\circ$:
\[
\fbi = r^+ p^+ o^\circ
\]

This gives Fizgig1.

Here $\fbo$ can either continue with $r^\circ$ or $r^-$.

\subsubsubsubsubsection{Quiddit}
Suppose $\fbo$ continue with $r^\circ$:
\[
\fbo = p^- r^\circ
\]

This forces
\[
\fbs = p^\circ
\]

This is Quiddit1.

But here $\fbb$ must eventually end at $s$ or $i$ after possibly twisting between the two. But then at least one of $\fbp$ and $\fbr$ will then escape to the outer perimeter and thus eventually absorbed.

This is a contradiction.


\subsubsubsubsubsection{Rudesby}
Back to the end of \ref{Fizgig}. Suppose $\fbo$ continued with $r^-$.

This is Rudesby.

Here $\fbs$ and $\fbo$ twist between $p$ and $r$ and eventually end there, since if either were to escape they would trace. Then as before,  $\fbb$ must eventually end at $s$ or $i$ after possibly twisting between the two. But then at least one of $\fbp$ and $\fbr$ will then escape to the outer perimeter and thus eventually absorbed.

This is a contradiction.

\subsubsubsubsection{Bawbee}
Back to the end of Xeric \ref{Xeric}. Suppose $\fbi$ continued with $o^+$.

This is Bawbee1.

Here $\fbo$, $\fbi$, and $\fbi$ are stuck in a three-way clockwise path mill between $p$ and $r$.

This is a contradiction.

\subsubsubsection{Bavin}
Back to the end of Fash \ref{Fash}. Suppose $\fbi$ continued with $s^*$. Then we must have
\[
\fbp = b^- o^\circ
\]
otherwise if the second letter of $\fbp$ were $o^-$ then it would immediately trace.

This is Bavin1.

Then in order for $\fbb$ to not trace, we must have the situation depicted in Bavin2, then continue with Bavin3. Here $\fbi$, $\fbr$, and $\fbb$ are stuck in a three-way clockwise path mill between $s$ and $i$.

This is a contradiction.

\subsubsection{Gley}
Back to the end of PPG \ref{PPG}. Suppose $\fbp$ continued with $s^*$. It cannot be $s^\circ$ as that would force $\fbb$ to begin with $s^+$ and immediately trace. So we must have
\[
\fbp = b^- s^- i^*...
\]

This is Gley1.

Here in order to prevent $\fbb$ from tracing, $\fbr$ must begin with $s^-$ to give us Gley2.

Here $\fbr$ can either continue with $p^*$ or $i^*$. Suppose it were the former. 

This is Gley3.

But here if $\fbr$ continued with $p^\circ$ or $p^+$ then $\fbs$ will immediately trace; whereas if $\fbr$ were to continue with $p^-$ then it immediately traces. This is a contradiction.

So suppose $\fbr$ continued with $i^*$. This is Gley4. Here $\fbs$ can either begin with $i^*$ or $p^*$. Suppose it were the former. Then this is Gley5.

Here if $\fbr$ were to continue with $i^\circ$ or $i^+$, then $\fbs$ will immediately trace; whereas if $\fbr$ were to continue with $i^-$ then it also immediately traces. This is a contradiction.

Thus suppose $\fbs$ began with $p^*$. Then $\fbo$ can begin with either $i^*$ or $p^*$. Suppose it were the former. This is Gley6.

Here if $\fbp$ continued into 1. or 2. then either $\fbr$ or $\fbo$ will be absorbed between $\fbp$ and $\fbi$. If $\fbp$ continued into 3 then it immediately is absorbed.

Thus it must be the case that $\fbo$ begins with $p$. This is Gley7. 

Here $\fbb$, $\fbp$, and $\fbr$ are stuck in a three-way clockwise path mill between $s$ and $i$.

This is a contradiction.

This concludes the case analysis for the assumptions under PPG \ref{PPG}.

\subsection{WVL}\label{WVL}
Suppose 
\[
\fbp = b^*...
\]
and 
\[
\fbi = b^*...
\]

This is WVL1.

Here $\fbi$ can end at $b^\circ$ immediately, or continue with $b^-$. Whereas $b^+$ is impossible a that would force $\fbp$ to continue with $b^+$ as well, and at least one of these will then immediately trace.

\subsubsection{Birl}\label{Birl}
Suppose
\[
\fbi = b^\circ
\]
Then we have the situation depicted in Birl1. 

Here $\fbb$ can either begin with $o^*$ or $s^*$. Suppose it were the former. Then $\fbr$ must begin with $o^-$ so as to not cause $\fbb$ to trace:
\[
\fbr = o^- i^*...
\]
Here it must continue with $i^-$. Otherwise, either  $\fbo$ or $\fbs$ will immediately trace.

This is Birl2.

Here we see that if $\fbr$ continued with $r$ then it trace, but if it continued with $b^-$ then it will circle back to the same position, so it cannot indefinitely continue with $b^-$. This is a contradiction.

Thus it must be the case that 
\[
\fbb = s^*...
\]
This forces $\fbo$ to begin with $i^-$. Otherwise, $\fbs$ will immediately trace.

This is Birl3.

Here $\fbs$ must begin with $i^-$, which means $\fbo$ must continue with $r^*$ or at least one of $\fbs$ or $\fbo$ will trace.

This is Birl4.

Here $\fbr$ can begin with $o^*$ or $s^-$. 

\subsubsubsection{Eisel}\label{Eisel}
Suppose 
\[
\fbr = o^*...
\]

Now $\fbs$ can either continue with $b^*$ or $o^*$. Suppose it were the former. Then we have Eisel1.

Here $\fbr$ must begin with $o^+$ and to the left of $\fbs$ or else $\fbs$ is immediately absorbed. But then $\fbr$ itself will be absorbed:
\[
\fbr = o^+ b^+ i^+ \Abso{\fbo}{\fbs}
\]

This is a contradiction. Thus $\fbs$ must continue with $r^*$.

This is Eisel2.

Here $\fbo$ can continue with $r^\circ$ or $r^+$, in which case $\fbs$ continues with $r^+$; or $\fbo$ can continue with $r^-$.

\subsubsubsubsection{Chouse}\label{Chouse}
Suppose $\fbo$ continued with $r^\circ$ or $r^+$, thus $\fbs$ continues with $r^+$.

Then $\fbs$ continues with $p^*...$ in order to not trace.

This is Chouse1.

Here $\fbb$ can begin with $s^-$ or $s^\circ$. If it were the former, then we have
\[
\fbb = s^- i^- r^*...
\]

which forces $\fbo$ to continue with $r^+$:
\[
\fbo = i^- \hat{r}^+ p^*...
\]
and thus
\[
\fbs = i^- r^+ \hat{p}^+...
\]

This is Chouse2.

Here $\fbb$, $\fbs$, and $\fbo$ are stuck in a three-way counterclockwise path mill between $p$ and $r$. This is a contradiction.

Thus we must have
\[
\fbb = s^\circ
\]

This gives us Chouse3. Here $\fbs$ and $\fbo$ twist $k$ times between $p$ and $r$ $k$ times to give us Candidate Family 10.

\subsubsubsubsection{Fardel}\label{Fardel}
Back to the end of Eisel \ref{Eisel}. Suppose $\fbo$ continued with $r^-$.

This is Fardel1.

Here $\fbp$ must continue with $i^-$. Otherwise, $\fbo$ will be forced to continue with $i^+$, in which case it will be absorbed between $\fbi$ and $\fbp$. Thus we have
\[
\fbp = b^+ \hat{i}^- r^*...
\]

This is Fardel2.

Here $\fbb$ can either begin with $s^-$, or begin and end with $s^\circ$. Suppose it were the former case. Then it continues:
\[
\fbb = s^- i^- r^- i^*...
\]
This is Fardel3.

Here $\fbs$ must continue with $r^-$ and to the right of $\fbp$. Otherwise $\fbp$ will be forced to continue with $r^+$ and then be absorbed between $\fbo$ and $\fbs$.

This leads to the situation depicted in Fardel4. Here $\fbp$ must twist between $r$ and $i$ and eventually end at either one; or it could end at $s$ after escaping. However, regardless which scenario, the remaining three strands, $\fbb$, $\fbo$ and $\fbs$ will be stuck in a three-way path mill. This is a contradiction.

Thus we must have
\[
\fbb = s^\circ
\]

This is Fardel5.

Here $\fbp$ must continue with $i^-$ in order to prevent $\fbo$ from being absorbed. 

This is Fardel6.

Here $\fbp$ twist $m$ times between $\fbi$ and $\fbr$ before ending at one of these as it cannot escape. Then the remaining $\fbo$ and $\fbs$ twist between $p$ and $i$ or $r$ $k$ times and end there.

This gives Candidate Family 11.

\subsubsubsection{Glaur}
Back to the end of Birl \ref{Birl}. Suppose 
\[
\fbr = s^-...
\]

But here $\fbr$ immediately traces.

This is a contradiction.

\subsubsection{Darg}\label{Darg}
Back to the end of WVL \ref{WVL}. Suppose $\fbi$ began with $b^-$. 

This is Darg1.

Here $\fbp$ can end immediately at $b^\circ$, or continue with $b^-$.

\subsubsubsection{Halse}\label{Halse}
Suppose 
\[
\fbp = b^\circ
\]

This is Halse1.

Here $\fbi$ can continue with $o^*$ or $s^*$.

\subsubsubsubsection{Kittle}\label{Kittle}
Suppose $\fbi$ were to continue with $o^*$. This is Kittle1.


Here it can be $o^\circ$, $o^-$, or $o^+$.

\subsubsubsubsubsection{Tantivy}
Suppose $\fbi$ continued with $o^\circ$:
\[
\fbi  = b^- o^\circ
\]

This is Tantivy1.

Here $\fbs$ can either begin with $i^*$ or $p^*$.

Suppose $\fbs$ begain with $i^*$. Then $\fbo$ also begins with $i^*$, and in fact
\[
\fbo = i^- r^*...
\]
This is Tantivy2.

Now consider $\fbr$. If it were to begin with $s^*$, then it must be $s^-$ or $\fbb$ will trace. But then the next letter will need to be $s^-$ and to the right of $\fbs$. Otherwise, $\fbs$ will immediately trace. But then $\fbr$ will trace:
\[
\fbr = s^- i^- r^\times
\]

Thus we must have
\[
\fbr = o^+ b^+ i^*...
\]
which forces 
\[
\fbs = i^- r^*...
\]

This is Tantivy3.

Here we claim that $\fbr$ must end here at $i^\circ$. Otherwise it will either trace:
\[
\fbr = o^+ b^+ \hat{i}^- r^\times
\]
or will be absorbed:
\[
\fbr = o^+ b^+ \hat{i}^+ b^- o^+ b^+ \Abso{\fbp}{\fbi}
\]

Thus we must have
\[
\fbr = o^+ b^+ \hat{i}^\circ
\]

This is Tantivy4.

Here $\fbb$ can either end immediately at $s^\circ$ or begin with $s^-$.

Suppose it were the latter, then we have
\[
\fbb = s^- i^- r^*...
\]

This is Tantivy5.

Here $\fbo$ and $\fbs$ continue with $r^+$ and to the left of $\fbb$, otherwise $\fbb$ will eventually trace. 

This is Tantivy6. 

But here $\fbo$, $\fbs$, and $\fbb$ are stuck in a three-way couterclockwise path mill between $p$ and $r$.

Thus it must be the former case, that
\[
\fbb = s^\circ
\]

This is Tantivy7.

Here $\fbs$ and $\fbo$ twist between $p$ and $r$ $k$ times and eventually end there. This gives us Candidate Family 12.

Now back to the situation depicted in Tantivy1. Suppose instead $\fbs$ began with $p^*$. Then $\fbo$ can either begin with $i^*$ or also begin with $p^*$. Suppose it were the former case:
\[
\fbo = i^*...
\]

Then we have Tantivy8.

Here we claim that $\fbr$ cannot begin with $s$. Indeed, if it did then it would have to be $s^-$, in which case the next letter must be $i^*$. If it were $i^\circ$ or $i^+$, then $\fbo$ will trace; if it were $i^-$ then $\fbr$ itself is immediately absorbed. Thus we must have
\[
\fbr = o^+ b^+ p^*...
\]
which forces 
\[
\fbs = p^- r^*...
\]

This is Tantivy9.

Here $\fbr$ must end. Otherwise, it either trace:
\[
\fbr = ... p^-r^\times
\]
or is absorbed:
\[
\fbr = ... p^+ b^- o^- b^+\Abso{\fbp}{\fbi}
\]

On the other hand, $\fbb$ and $\fbo$ both must also end, as there is no possible way for them to twist without at least one of them being immediately absorbed. Then $\fbs$ must also end.

This gives Tantivy10, which is Candidate 13.

Now suppose it were the case that $\fbo$ began with $p^*$. This is Tantivy11. 

Here we claim that $\fbr$ cannot begin with $o^+$. Indeed, if it did, then it must be
\[
\fbr = o^+ b^+ p^*...
\]
which forces
\[
\fbs = p^- r^*...
\]

This is Tantivy12. But here $\fbr$ must end, otherwise it would either immediately trace if continuing with $p^-$, or eventually absorbed if continuing with $p^+$. But then $\fbs$ must continue with $r^+$ and to the left of $\fbo$. Otherwise $\fbo$ will be absorbed. Now $\fbs$ will trace:
\[
\fbs = p^- \hat{r}^+ p^+ s^\times
\]

Thus we must have
\[
\fbr = s^-...
\]
and the next letter must be $i^*$.

This is Tantivy13.

Here suppose further that $\fbo$ continued with $r^+$:
\[
\fbo = p^- \hat{r}^+p^+ s^*...
\]
which forces
\[
\fbb = s^-...
\]
if the next letter were $r^-$, then we have Tantitvy14, where $\fbs$ must begin with $p^-$:
\[
\fbs = p^- r^+ p^+ s^\times
\]

So this cannot occur. Suppose instead the next letter for $\fbb$ were $i^*$. Then we have Tantivy15, where $\fbo$, $\fbr$, and $\fbb$ are stuck in a three-way clockwise path mill between $s$ and $i$.

Back to the situation depicted in Tantivy13. Suppose $\fbo$ were to end there:
\[
\fbo = p^- r^\circ
\]

This forces $\fbs$ to end immediately otherwise it will eventually trace. This gives Tantivy16. Here $\fbb$ and $\fbr$ twist between $s$ and $i$ $k$ times and eventually end there, giving a candidate that will be subsumed in Candidate Family 15 to follow.

Back to the situation depicted in Tantivy13. Suppose $\fbo$ continued with $r^-$:
\[
\fbo = p^- r^-p^*...
\]

This forces
\[
\fbs = p^- r^*...
\]

This is Tantivy17. 

Here $\fbo$ and $\fbs$ twist $k$ times between $p$ and $r$ times and end there; while $\fbr$ twist $m$ times between $s$ and $i$ $k$ times and end there. This gives Candidate Family 14.


\subsubsubsubsubsection{Wherret}
Back to the end of Kittle \ref{Kittle}. Suppose $\fbi$ were to continue with $o^-$. Then $\fbo$ and $\fbs$ must begin with $p^*$, and $\fbr$ must begin with
\[
\fbr = o^- p^*...
\]

This is Wherre1.

Here $\fbr$ and $\fbi$ must continue with $p^-$ and to the right of $\fbo$ and $\fbs$. Otherwise, $\fbs$ would begin with $p^+$, in which case it would immediately trace. Thus we have
\[
\fbr = s^- \hat{p}^- b^- o^- i^*...
\]
and 
\[
\fbi = b^- o^- \hat{p}^- r^*...
\]

This is Wherret2.

Here $\fbs$ must begin with $p^+$, otherwise $\fbo$ would trace. Thus we have
\[
\fbs = p^+ o^*...
\]
which forces
\[
\fbb = o^- p^- b^\times
\]

This is a contradiction.

\subsubsubsubsubsection{Xiphias}\label{Xiphias}
Back to the end of Kittle \ref{Kittle}. Suppose $\fbi$ were to continue with $o^+$.

This is Xiphias1.

Here $\fbr$ can begin with $o^+$ or $s^*$.

Suppose it were the former:
\[
\fbr = o^+b^+ p^*...
\]

Here $\fbs$ can either begin with $i^+$ or $p^*$. Suppose it were the former, then it forces $\fbb$ to begin with $o^+$, and then immediately trace. As shown in Xiphias2. Thus $\fbs$ must begin with $p^*$. Now $\fbo$ can begin with $i^*$ or $p^*$. Suppose it were the former. Then it must in fact be 
\[
\fbo = i^\circ
\]
Otherwise it would cause $\fbb$ to be absorbed. This is Xiphias3.

Here $\fbs$ continues to give us Xiphias4.

Here $\fbr$ cannot continue with $p^\circ$. This would imply that $\fbi$ continus with $p^- o^*...$, and then forcing $\fbs$ to continue with $o^-...$, as shown in Xiphias5. Here $\fbs$ will eventually be absorbed in between $\fbp$ and $\fbi$. On the other hand, it is clear that $\fbr$ cannot continue with $p^-$ as that would mean it either immediately trace, or escape to the outer perimeter and absorbed between itself and $\fbb$. So it must continue with $p^+$, as shown in Xiphias6.

Here we claim that $\fbs$ must end at $o^\circ$. Indeed, suppose it did not; then it continues with $o^-$, to give us the situation depicted in Xiphias7. But now $\fbr$, $\fbs$, and $\fbi$ are stuck in a three-way counterclockwise path mill between $o$ and $p$. Thus we must have
\[
\fbs = p^- o^\circ
\]
to give us Xiphias8.

Here $\fbr$ and $\fbi$ can twist between $p$ and $o$ (traversing counterclockwise) $k$ times and must end at either of those locations. This gives us Candidate Family 15.

Suppose now that $\fbo$ began with $p^*$. 

This is Xiphias9

Here $\fbo$ must continue with $p^-$ and to the right of $\fbr$. Otherwise, $\fbs$ will be forced to begin with $i^+$ and immediately trace. This then also forces $\fbs$ to begin with $p^-$. Otherwise, $\fbi$ and $\fbr$ will be absorbed between $\fbo$ and $\fbs$.

This is Xiphias10

Here $\fbo$ must continue with $r^-$ and to the right of $\fbs$. Otherwise, $\fbs$ will be forced to continue with $r^+$ and trace:
\[
\fbs = p^- \hat{r}^+ p^+ s^\times
\]

Thus we have Xiphia11

Here if $\fbr$ continued with $p^\circ$ or $p^+$ then $\fbo$ will trace:
\[
\fbo = ...p^+ b^- o^\times
\]
But if $\fbr$ continued with $p^-$ it immediately traces.

This is a contradiction.

Back to the situation depicted in Xiphias1. Suppose$\fbr$ began with  $s^*$.

This is Xiphias12

Here $\fbo$ can begin with $i$ or $p$. Suppose it were the former, then it must be $s^-$ otherwise $\fbs$ will immediately trace. But now $\fbo$ traces:
\[
\fbo = i^- b^- o^\times
\]

Thus we must have
\[
\fbo = p^- r^*...
\]

This gives Xiphias13

But here $\fbi$, $\fbo$, and $\fbs$ are stuck in a three-way clockwise path mill between $p$ and $r$. This is a contradiction.


\subsubsubsubsection{Nidder}
Back to the end of Halse \ref{Halse}. Suppose $\fbi$ continued with $s^*$:
\[
\fbi b^- s^*...
\] 

This is Nidder1.

Here $\fbr$ and $\fbi$ must continue with $s^-$ and to the right of $\fbb$. Otherwise, $\fbb$ would begin with $s^+$ and immediately trace.

This is Nidder2.

Here $\fbo$ and $\fbr$ must continue with  $p^-$ and to the right of $\fbs$. Otherwise, $\fbs$ would begin with $p^+$ and then immediately trace.

This gives us Nidder3.

Here notice that $\fbr$ circles the outer perimeter clockwise some number of times, and eventually it must go up to $i^*$, as the only other choice is $r$, but then it would trace.

Nidder4 shows an example. While $\fbr$ circles, $\fbi$ will follow. Once $\fbr$ goes up to $i^*$, then $\fbi$ must then go to $r$, otherwise it would follow $r$ to $i^*$ and trace.

Now $\fbo$ cannot also go up to $i^*$, since in that case, either $\fbr$ continued with $i^\circ$ or $i^+$, forcing $\fbo$ to continue with $i^+$, then it will be eventually absorbed; but then if $\fbr$ went $i^-$, then $\fbr$ will eventually be absorbed as well. This can be seen in Nidder4.

Thus we must be in the situation depicted in Nidder5. But here if $\fbo$ continued with $r^\circ$ or $r^+$, forcing $\fbi$ to continue with $r^+$, then $\fbi$ will be absorbed:
\[
\fbi = ...r^+ p^+ s^+ b^+ p^+ s^+ b^+ p^+ s^+ \Abso{\fbr}{\fbb}
\]

On the other hand, if $\fbo$ continued with $o^-$, then it would also be absorbed:
\[
\fbo = ...r^- p^+ s^+ b^+ p^+ s^+ b^+ \Abso{\fbp}{\fbi}
\]

This is a contradiction.

\subsubsection{Busk}\label{Busk}
Back to the end of WVL \ref{WVL}. Suppose $\fbi$ and $\fbp$ both began with $b^-$
\[
\fbi = b^-...
\]
and
\[
\fbp = b^-...
\]

This is Busk1.

Suppose $\fbi$ continued with $s*$. Then $\fbi$ will trace, as depicted in Busk2.

So $\fbi$ must continue with $o^*$. On the other hand $\fbp$ can continue with $o^*$ as well or continue with $s^*$ in stead. Suppose it were the former.

This is Busk3.

Here $\fbb$ can begin with $o^*$ or $s^*$. Suppose it were the former case. But now $\fbi$ traces, as shown in Busk4.

So $\fbb$ must begin with $s^*$, and $\fbr$ can also begin with $s^*$ or begin with $o^*$.

\subsubsubsection{Kieve}\label{Kieve}
Suppose $\fbr$ began with $s^*$. Then it must be $s^-$, otherwise $\fbb$ would trace. The second letter can be $p^*$ or $i^*$. Suppose it were the former. Then this forces
\[
\fbs = p^*...
\]
and
\[
\fbo = p^- r^*...
\]

This is Kieve1.

But now $\fbr$ trace.

So it must be 
\[
\fbr = s^- i^*...
\]

Now $\fbs$ can begin with  $i^*$ or $p^*$. Suppose it were the former. Then this forces
\[
\fbo = i^-...
\]

and the next letter cannot be $r$ or it would force $\fbr$ to trace. Thus we must have the situation depicted in Kieve2.

But now $\fbo$ will trace:
\[
\fbo = s^- b^- o^\times
\]

Thus $\fbs$ must begin with $p^*$. Then $\fbo$ can either begin with $i^+$, or also begin with $p^*$. Suppose it were the former. Then it must do so to the left of the incoming $\fbr$, otherwise $\fbr$ would be immediately absorbed. This forces the situation depicted in Kieve3. Here $\fbo$, $\fbi$, and $\fbp$ are stuck in a three-way path mill between $b$ and $o$. This is a contradiction.

Thus we must have
\[
\fbo = p^- r^*...
\]

This gives Kieve4.

Here three are three twisting pairs: $\fbi$ and $\fbp$ twist $k$ times between $b$ and $o$; $\fbb$ and $\fbr$ twist $m$ times between $s$ and $i$; and $\fbs$ and $\fbo$ twist $j$ times between $p$ and $r$. Each of the image paths of these pairs cannot escape and enter into another pair, as that would immediately create a three-way path mill. This gives us Candidate Family 16.

\subsubsubsection{Oorie}
Back to the end of Busk \ref{Busk}. Suppose and $\fbr$ began with $o^*$. Then it must be 
\[
\fbr = b^+...
\]
otherwise $\fbi$ and $\fbp$ will immediately trace.  This is Oorie1.

The next letter cannot be $i$, as this would force the situation depicted in Oorie2, where $\fbo$ must begin with $i^\circ$ or $i^+$ or it will eventually trace; but then this forces $\fbs$ to begin with $i^+$ and immediately trace.

Thus we must have
\[
\fbr = o^+ b^+ p^*...
\]

Now again $\fbs$ cannot begin with $i^*$ for the reason as before, thus it must begin with $p^*$, and continue:
\[
\fbs = p^- r^*...
\]

Now $\fbo$ can begin with $i^*$ or $p^*$ as well. Suppose it were the former. 

This is Oorie3.

Then it must be the case that 
\[
\fbb = s^\circ
\]
and
\[
\fbo = i^\circ
\]

Indeed, otherwise, one of $\fbb$ and $\fbo$ would force the other to be immediately absorbed.

This is Oorie4.

Now we claim that $\fbs$ must end here at $r^\circ$. Suppose that were not the case, then it must continue with $r^-$, then it continues:
\[
\fbs = ...r^- i^+ b^- o^*...
\]

This is Oorie5.

Here $\fbp$ will twist between $o$ and $b$ and eventually end there, as escaping would force it to trace. But once it ends, the remaining $\fbs$, $\fbi$, and $\fbr$ will be stuck in a three-way path mill. 

For example, Oorie6 shows a scenario where it is clear that they are stuck in a counterclockwise path mill. The other situations will be the same except with more twisting of $\fbp$ before it ends.

Thus we must have
\[
\fbs = p^- r^\circ
\]

This is Oorie7.

Here again $\fbp$ twist between $b$ and $o$ some number of times, say $k$ times. Then $\fbi$ and $\fbr$ twist $m$ times between $p$ and $o$ or $b$, $k$ times and end there.

This is Candidate Family 17.

Oorie8 shows an example.


Now suppose  $\fbo$ began with $p^*$. Then it must be $p^-$ as to not make $\fbs$ trace.

This forces the situation depicted in Oorie9.

Here $\fbr$ must continue with $p^\circ$ or $p^+$, or it will immediately trace. But then this forces $\fbo$ to continue with $p^+$, and it will trace:
\[
\fbo = p^- r^- \hat{p}^+ b^- o^\times
\]

This is a contradiction.


\subsubsubsection{Quank}
Back to the middle part of Busk \ref{Busk}. Suppose
\[
\fbi = b^- o^*...
\]
and 
\[
\fbp = b^- s^*...
\]

This is Quank1.

Here $\fbs$ can begin with $i^*$ or $p^*$. Suppose it were the former. Then this forces
\[
\fbo = i^- r^*...
\]

which then forces
\[
\fbp = b^- s^- i^- r^*...
\]

This is Quank2.

But now $\fbr$ will trace.

Thus we must have
\[
\fbs = p^*...
\]

and $\fbo$ can either begin with $i^*$ or $p^*$. Suppose it were the former. Then it continues and is absorbed:
\[
\fbo = i^+ s^+ b^+ \Abso{\fbp}{\fbi} 
\]

Thus we must have
\[
\fbo = p^- r^*...
\]

This is Quank5.

Here $\fbr$ must continue with $i^*...$, but then this forces $\fbp$ to continue with $i^-$ and to the right of the incoming $\fbi$.  But now $\fbb$, $\fbr$, and $\fbp$ are stuck in a three-way clockwise path mill between $s$ and $i$, shown in Quank5.

This finishes the case analysis for the Candelabra.